%% ============================================================
%% 分片 13：原文第 5 节「Topological Field Theory In Two Dimensions」（原文 glshort.tex 起于第 3264 行）
%% 由 tools/harvmac2tex.py 从 arXiv 源码（hep-th/0604151v3）自动生成。
%% 只译英文散文；公式、编号、\label、\cite 一律不改。
%% ============================================================

\subsection{The Role Of Generalized Complex Geometry}

\label{sec:gencpx}

Generalized complex geometry is a natural framework for describing
the topological field theories that can be constructed by twisting
two-dimensional sigma-models, such as the sigma-model with target
$\MH$.  Our goal here is to use generalized complex geometry to
clarify a few questions from a two-dimensional point of view. What
is the family $\Bbb{CP}^1_g$ and why is it independent of the
complex structure of $C$?  How, from a two-dimensional point of
view, can we understand the canonical parameter $\Psi$ introduced
in section 第~\ref{sec:canonpar}~节?   And what geometry of $\MH$ is really
needed in these constructions?

 Let $TX$ and $T^*X$ be the tangent
and cotangent bundles of a manifold $X$. We let $\hat T=TX\oplus
T^*X$ and we write a section of $\hat T$ as $\left(\begin{matrix}v\\
\xi\\\end{matrix}\right)$, where $v$ is a section of $TX$ and $\xi$ is a
section of $T^*X$.  $\hat T$ has a natural indefinite metric, in
which $TX$ and $T^*X$ are both null and the inner product of $TX$
and $T^*X$ is the natural one, in which the inner product  of
$\left(\begin{matrix}v\\ \xi\\\end{matrix}\right)$ and $\left(\begin{matrix}v'\\
\xi'\\\end{matrix}\right)$ is $v^i\xi'_i+v'{}^i\xi_i$.

  
 

 A generalized
almost complex structure on a manifold $X$ is a linear
transformation ${\cal I}$ of $\widehat T=TX\oplus T^*X$ that
preserves the metric and obeys ${\cal I}^2=-1$.  If a certain
integrability condition is obeyed, it is called a generalized
complex structure \cite{ref69}.  (For more detail, see Gualtieri's
thesis \cite{ref67}, as well as \cite{ref71,ref72,ref73,ref74,ref68} for applications
to sigma-models and \cite{ref75,ref76,ref77} for applications to
supergravity.) One basic example of a generalized complex
structure is % 式 (5.23)
\setcounter{equation}{22}
\begin{equation}\label{eq:momo}
{\cal I}_{\hat J}=\left(\begin{matrix} {\hat J} &
0 \\ 0 & -{\hat J}\,^t\\\end{matrix}\right),
\end{equation}
 where ${\hat J}:TX\to TX$ is
an ordinary complex structure, and ${\hat J}\,^t:T^*X\to T^*X$ is
its transpose.  If $\omega$ is a symplectic structure, then a
second basic example of generalized complex structure is given by
% 式 (5.24)
\setcounter{equation}{23}
\begin{equation}\label{eq:zomo}
{\cal I}_\omega=\left(\begin{matrix}0 & -\omega^{-1}\\ \omega
& 0\\\end{matrix}\right).
\end{equation}
 Here we regard $\omega$ as a linear map from $TX$
to $T^*X$; $\omega^{-1}$ is the inverse map from $T^*X$ to $TX$.
The topological field theory associated with ${\cal I}_{\hat J}$
is the $B$-model in complex structure ${\hat J}$, and the one
associated with ${\cal I}_\omega$ is the $A$-model with symplectic
structure $\omega$.\footnote{The $A$-model is most commonly considered
on a Kahler manifold, and then $\omega$ is the Kahler form.
Because the $A$-model makes sense on symplectic manifolds more
generally, we simply refer to $\omega$ as the symplectic form.}

In each of these cases, the $B$-field vanishes.  A $B$-field can
be turned on as follows.  For any closed two-form $B_0$, let
% 式 (5.25)
\setcounter{equation}{24}
\begin{equation}\label{eq:hyt}
{\cal M}(B_0)=\left(\begin{matrix}1 & 0 \\
                                   B_0 & 1 \\\end{matrix}\right).
\end{equation}

Then if ${\cal I}$ is an integrable generalized complex structure,
so is % 式 (5.26)
\setcounter{equation}{25}
\begin{equation}\label{eq:nyt}
{\cal I}(B_0)={\cal M}(B_0){\cal I}{\cal
M}(B_0)^{-1}.
\end{equation}
 The transformation ${\cal I}\to {\cal I}(B_0)$ has
the effect of shifting the $B$-field by $B_0$.  So in particular
the topological field theory derived from ${\cal I}_{\hat J}(B_0)$
is the $B$-model with complex structure ${\hat J}$ and $B$-field
$B_0$, and ${\cal I}_\omega(B_0)$ is similarly related to the
$A$-model with symplectic structure $\omega$ and $B$-field $B_0$.
Conjugation by ${\cal M}(B_0)$ is called the $B$-field transform.

It is shown in chapter 6 of \cite{ref67}\ that the conditions \cite{ref66}\
for a sigma model to have $(2,2)$ supersymmetry are equivalent to the existence of a pair
of generalized complex structures obeying a certain compatibility
condition.  If the $(2,2)$ model has anomaly-free $R$-symmetries, then it can be twisted
in two ways to make a topological field theory.  It is believed that, in general, each of these
topological field theories is determined by one of the two generalized complex structures and independent
of the second.
(For example, if the target is a Kahler manifold and the $B$-field is flat, then the two
generalized complex structures are ${\cal I}_{\hat J}$ and ${\cal I}_\omega$, as described above;
the two topological field theories are the $B$-model, which is determined by the first, and the
$A$-model, which is determined by the second.)  Arguments supporting this claim can be found in
\cite{ref72,ref68}, though a complete proof is not yet known.

In the case relevant to us that the $(2,2)$ model
comes from a hyper-Kahler metric $g$ with a pair of points
$J_+,J_-\in \Bbb{CP}^1_h$, there is a slight simplification in the formulas
of \cite{ref67},
because the $B$-field can be taken to vanish (and restored later
by a $B$-field transform).  We let $\omega_\pm$ be the two
symplectic structures $\omega_\pm = gJ_\pm$.  One of the two
generalized complex structures determined by the pair $J_+$, $J_-$
with the hyper-Kahler metric $g$ is then, according to eqn. 6.3 of
\cite{ref67}, % 式 (5.27)
\setcounter{equation}{26}
\begin{equation}\label{eq:ugg}
{\cal J}={1\over 2}\left(\begin{matrix}J_++J_- &
-(\omega_+^{-1}-\omega_-^{-1})\\ \omega_+-\omega_- &
-(J_+^t+J_-^t)\\\end{matrix}\right).
\end{equation}
 The second one, not relevant for us,
is obtained by reversing the sign of $J_-$ and $\omega_-$.

Using \eqref{eq:wyrre}\ and $(w_+,w_-)=(-t,t^{-1})$, we have % 式 (5.28)
\setcounter{equation}{27}
\begin{equation}\label{eq:ytro}
\begin{aligned} J_+ & = {1-\bar t t\over 1+\bar t t}I-{i(t-\bar t)\over
1+\bar t t}J-{t+\bar t\over 1+\bar t t}K  \\ J_- & = -{1-\bar
tt\over 1+\bar t t}I-{i(t-\bar t)\over 1+\bar t t}J+{t+\bar t\over
1+\bar t t}K  .\\\end{aligned}
\end{equation}
 The associated symplectic structures are
 % 式 (5.29)
\setcounter{equation}{28}
\begin{equation}\label{eq:bytro}
\begin{aligned} \omega_+ & ={\rm Im}\,\tau\left( {1-\bar t t\over 1+\bar
t t}\omega_I-{i(t-\bar t)\over 1+\bar t t}\omega_J-{t+\bar t\over
1+\bar t t}\omega_K\right) \\ \omega_- & = {\rm
Im}\,\tau\left(-{1-\bar tt\over 1+\bar t t}\omega_I-{i(t-\bar
t)\over 1+\bar t t}\omega_J+{t+\bar t\over 1+\bar t
t}\omega_K\right) .\\\end{aligned}
\end{equation}
 (The factor of ${\rm Im}\,\tau$ was
obtained in \eqref{eq:inog}.) So we compute that the generalized complex
structure determined by this data is % 式 (5.30)
\setcounter{equation}{29}
\begin{equation}\label{eq:islo}
{\cal I}_t={1\over
1+\bar t t}\left(\begin{matrix}-i(t-\bar t)J & -({\rm
Im}\,\tau)^{-1}((1-\bar t t)\omega_I^{-1}-(t+\bar
t)\omega_K^{-1})\\ {\rm Im}\,\tau ((1-\bar t t)\omega_I-(t+\bar
t)\omega_K)& i(t-\bar t)J^t\\\end{matrix}\right).
\end{equation}
 (To get this formula, it
helps to know that on a hyper-Kahler manifold, if $a^2+b^2+c^2=1$,
then
$(a\omega_I+b\omega_J+c\omega_K)^{-1}=a\omega_I^{-1}+b\omega_J^{-1}+c\omega_K^{-1}$.)

The first important observation is that $I, K$, and $\omega_J$
have disappeared.  ${\cal I}_t$ depends only on $J$, $\omega_K$,
and $\omega_I$, or equivalently on $J$ and the holomorphic
two-form $\Omega_J=\omega_K+i\omega_I$ (and its complex
conjugate). But this data, as we noted in section 第~\ref{sec:mhhyper}~节, is
independent of the metric of $C$. This explains, from the
two-dimensional point of view, why the family of topological field
theories that we get by dimensional reduction from four dimensions
does not depend on the metric.

An analogous family of generalized complex structures can be
constructed on any hyper-Kahler manifold $X$ and is described in
section 4.6 of \cite{ref67}.  As is explained there (Proposition
4.34), the generic element in such a family is a $B$-field
transform of a generalized complex structure derived from a
symplectic structure. Indeed, a small calculation shows that
% 式 (5.31)
\setcounter{equation}{30}
\begin{equation}\label{eq:huffy}
{\cal I}_t={\cal M}(B_0)\left(\begin{matrix}0 &
-\omega_0^{-1}\\ \omega_0 & 0\\\end{matrix}\right){\cal M}(B_0)^{-1},
\end{equation}

where % 式 (5.32)
\setcounter{equation}{31}
\begin{equation}\label{eq:nutty}
\begin{aligned}\omega_0 & = {\rm Im}\,\tau{1-\bar
t{}^2t^2\over (1+t^2)(1+\bar t{}^2)}\left(\omega_I-{t+\bar t\over
1-\bar t t}\omega_K\right)\\
                      B_0 &= -{\rm Im}\,\tau\left.{i(t^2-\bar t{}^2)\over(1+t^2)(1+\bar
                      t{}^2)}\right.
                      \left(\omega_I+{1-\bar t t\over t+\bar t}\omega_K\right).\\\end{aligned}
\end{equation}

Therefore, for $t\neq \pm i$ the TQFT derived from the generalized
complex structure ${\cal I}_t$ is equivalent to an $A$-model with
symplectic structure $\omega_0$ and $B$-field $B_0$.  This fact
will enable us to understand from a two-dimensional point of view
the canonical parameter $\Psi$ introduced in section 第~\ref{sec:canonpar}~节.

In general, the $A$-model with symplectic form $\omega_0$ and
$B$-field $B_0$ depends only on the cohomology class
$[B_0+i\omega_0]$.  (On a Kahler manifold, this is called the
complexified Kahler class.)  In the present problem, since the
cohomology class of $\omega_K$ vanishes (as we explained in
section 第~\ref{sec:mhhyper}~节), we have % 式 (5.33)
\setcounter{equation}{32}
\begin{equation}\label{eq:gutty}
[B_0+i\omega_0]=-i\,{\rm
Im}\,\tau\left(t-t^{-1}\over t+t^{-1}\right) [\omega_I].
\end{equation}
 Thus,
when the four-dimensional $\theta$ angle vanishes (we took it to
vanish by starting with ${\cal I}_t$ rather than a $B$-field
transform thereof), the model depends on ${\Im}\tau$ and $t$ only
in the combination that appears in \eqref{eq:gutty}.

The four-dimensional $\theta$ angle induces, as we explained at
the end of section 第~\ref{sec:mhhyper}~节, an additional contribution
$B'=-(\theta/2\pi)\omega_I=-({\rm Re}\,\tau)\,\omega_I$ to the
two-dimensional $B$-field. This can be incorporated in the
generalized complex structure simply by conjugating with ${\cal
M}(B')$.  The resulting model depends on the cohomology class
$[B_0+B'+i\omega_0]$, which is $-[\omega_I]$ times
% [WARN] 公式 \eqn\eqref{eq:loopy} 无编号或重复
\begin{equation*}
\tag*{(5.34)}
\Psi={\rm Re}\,\tau +i\,{\rm Im}\,\tau
\left(t-t^{-1}\over t+t^{-1}\right).
\end{equation*}
 This gives a two-dimensional
interpretation of why the model depends on $t$ and $\tau$ only in
the combination $\Psi$.

\par\medskip{\it A Few Loose Ends}

Finally, let us wrap up a few loose ends.

First of all, if $t$ is real, the TQFT determined by complex
structures $J_+(t),J_-(t)$ is an $A$-model to begin with, and the
above argument showing that $\Psi$ is the only relevant parameter
did not really require generalized complex geometry.   We could have based our
reasoning on the real case together with holomorphy in $t$.     We chose not to do this since
 we think that generalized
complex geometry is a natural framework for understanding this problem, even though there are
some technical gaps in the current understanding of it.

Second, in the above analysis, we used the fact that the $A$-model
only depends on the cohomology class of $B+i\omega$.  This is
proved by writing the action as $\{Q,V\}+\int_\Sigma
\Phi^*(\omega-iB)$, where  $\Phi:\Sigma\to X$ is the sigma-model
map.

In the above derivation, the $B$-field is not just of type $(1,1)$
(in a complex structure in which the symplectic form is Kahler);
it also has components of type $(2,0)\oplus (0,2)$. In fact, the
$A$-model is sensitive to all components of the $B$-field,
including the part of type $(2,0)\oplus (0,2)$, but this point may
require some clarification.

If $\Phi:\Sigma\to X$ is a holomorphic map, and $B$ is a form on
$X$ of type $(2,0)\oplus (0,2)$, then $\Phi^*(B)=0$.  This makes
the $(2,0)\oplus (0,2)$ part of $B$ appear irrelevant, if one
interprets the $A$-model purely as a mechanism for computing
correlation functions by summing over holomorphic maps. But if one
considers branes (as we most definitely will do to understand the
geometric Langlands program), one immediately sees that all parts
of the Hodge decomposition of $B$ are relevant.  A Lagrangian
submanifold $N\subset X$ endowed with a Chan-Paton line bundle
${\cal L}$ of curvature $F$ gives an $A$-brane if $F+B|_N=0$. This
condition is certainly sensitive to the $(2,0)\oplus (0,2)$ part
of $B$.

Finally, our analysis showing that $\Psi$ is the only relevant
parameter is really not valid at $t=\pm i$, because of poles in
the formulas.  At these values of $t$, the model is actually not
the $B$-field transform of an $A$-model; it is a $B$-model in
complex structure $\pm J$. To complete our analysis for these
values of $t$, we will argue directly using the $B$-model.

At $t=\pm i$, we have $\Psi=\infty$, independent of $\tau$, so to
complete the demonstration that $\Psi$ is the only relevant
parameter, we must show that $\tau$ is completely irrelevant at
$t=\pm i$. ${\Im}\,\tau$ controls the Kahler class of $\MH$, and
this is certainly irrelevant in the $B$-model.  Varying ${\rm
Re}\,\tau$ adds to the $B$-field a multiple of $\omega_I$. To show
that this term is irrelevant in the $B$-model with complex
structure $J$, we write $\omega_I=-i\Omega_J+i\omega_K$. The
contribution from $\omega_K$ is irrelevant because the form
$\omega_K$ is exact. The contribution from $\Omega_J$ is
irrelevant because $\Omega_J$ is a form of type $(2,0)$.  But in
the $B$-model, the $B$-field contribution $-i\int_\Sigma\Phi^*(B)$
can be written as $\{Q,\dots\}$ if $B$ is of type $(1,1)$ or
$(2,0)$.  (By contrast, a $(0,2)$ component of the $B$-field does
affect the category of $B$-branes \cite{ref71}.) Of course, in complex
structure $-J$, we make the same argument, starting with
$\omega_I=i\bar\Omega_J-i\omega_K$.

\subsection{Some Specializations}

\label{sec:somespec}

In the rest of this paper except section 第~\ref{sec:genotics}~节, we focus
primarily on the distinguished values $\Psi=0$ and $\Psi=\infty$,
and the duality % 式 (5.35)
\setcounter{equation}{34}
\begin{equation}\label{eq:rtu}
S=\left(\begin{matrix}0 & 1/\sqrt {n_{\frak
g}}\\ -\sqrt {n_{\frak g}}& 0\\\end{matrix}\right)
\end{equation}
 that maps $\Psi\to
-1/n_{\frak g}\Psi$ and hence exchanges $\Psi=0$ and
$\Psi=\infty$. We recall that $n_{\frak g}=1$ for simply-laced
$G$. When we speak of $S$-duality, we will be a little imprecise
about whether we mean the full duality group $\Gamma$ generated by
$S$ and by % 式 (5.36)
\setcounter{equation}{35}
\begin{equation}\label{eq:ntu}
T=\left(\begin{matrix}1 & 1\\ 0& 1\\\end{matrix}\right),
\end{equation}
 or
just the subgroup generated by $S$. In fact,  a complete story
requires including also $T$, but $S$ is the main actor in the
geometric Langlands program.

To get $\Psi=\infty$, we take $t=i$ and thus $(w_+,w_-)=(-i,-i)$.  The resulting model
is the $B$-model in complex structure $J$.  This
statement is unaffected by the choice of $\tau$, but it is
convenient to take $\tau$ imaginary.

Applying the transformation $S$, we find that $\tau$ remains
imaginary, and $t$ is mapped to $t'=-t ({ \tau}/|\tau|)=-it=1$.
This leads to $(w_+,w_-)=(-1,1)$, and hence to the $A$-model in
complex structure $K$.  Thus $S$-duality implies that the
$B$-model for $\MH(G,C)$ in complex structure $J$ is equivalent to
the $A$-model for $\MH(\LG,C)$ in complex structure $K$. If we
combine $S$ with a transformation ${\EUI}:\varphi\to i\varphi$
that rotates complex structure $K$ back to complex structure $J$,
we get a mirror symmetry that maps the $B$-model in complex
structure $J$ to the $A$-model in complex structure $J$.  So the
complex manifolds $\MH(G,C)$ and $\MH(\LG,C)$ are a mirror pair in
complex structure $J$.

The automorphism ${\EUI}$ that we used here is not quite natural,
however, in the sense that it moves us out of the four-dimensional
family of TQFT's parametrized by $\Bbb{CP}^1_g$. It has to move us
out of this family, since the $A$-model in complex structure $J$
is not a member of this family!

Additionally, we already noted in discussing \eqref{eq:yfog}\ that $S$ (and
the whole duality group) leaves invariant the $B$-model in complex
structure $I$, and likewise the $A$-model in that complex
structure.  And finally, since $S$ exchanges the $B$-model in
complex structure $J$ with the $A$-model in complex structure $K$,
we may guess that it likewise exchanges the $A$-model in complex
structure $J$ with the $B$-model in complex structure $K$. In
fact, the automorphism ${\EUI}$ relates these two statements, but
instead of arguing this way, let us use eqn. \eqref{eq:yfog}.  The $A$-model
in complex structure $J$ corresponds to $(w_+,w_-)=(-i,i)$.  This
is mapped by $S$ to $(w_+,w_-)=(1,1)$, which indeed corresponds to
the $B$-model in complex structure $K$. We summarize these
statements in a table.

\par\medskip

\begin{table}[htbp]
\centering
\begin{tabular}{ll}
\toprule
Model & $S$-Dual \\
\midrule
$I_B$ & $I_B$ \\
$I_A$ & $I_A$ \\
$J_B$ & $K_A$ \\
$J_A$ & $K_B$ \\
$K_B$ & $J_A$ \\
$K_A$ & $J_B$ \\
\bottomrule
\end{tabular}
\caption{A model and its $S$-dual.  $I_B$, for example, represents the $B$-model in complex structure $I$, which is its own $S$-dual.（原文 Table~2）}
\label{tab:sdual-models}
\end{table}
